\chapter{Data structures}

\index{data structure}

A \key{data structure} is a way to store
data in the memory of a computer.
It is important to choose an appropriate
data structure for a problem,
because each data structure has its own
advantages and disadvantages.
The crucial question is: which operations
are efficient in the chosen data structure?

This chapter introduces the most important
data structures in the Rust standard library.
It is a good idea to use the standard library
whenever possible,
because it will save a lot of time.
Later in the book we will learn about more sophisticated
data structures that are not available
in the standard library.

\section{Dynamic arrays}

\index{dynamic array}
\index{vector}

A \key{dynamic array} is an array whose
size can be changed during the execution
of the program.
The most popular dynamic array in Rust is
the \texttt{Vec} type,
which can be used almost like an ordinary array.

The following code creates an empty vector and
adds three elements to it:

\begin{lstlisting}
let mut v = Vec::new();
v.push(3); // [3]
v.push(2); // [3,2]
v.push(5); // [3,2,5]
\end{lstlisting}

After this, the elements can be accessed like in an ordinary array:

\begin{lstlisting}
println!("{}", v[0]); // 3
println!("{}", v[1]); // 2
println!("{}", v[2]); // 5
\end{lstlisting}

The method \texttt{len} returns the number of elements in the vector.
The following code iterates through
the vector and prints all elements in it:

\begin{lstlisting}
for i in 0..v.len() {
    println!("{}", v[i]);
}
\end{lstlisting}

\begin{samepage}
A shorter way to iterate through a vector is as follows:

\begin{lstlisting}
for x in &v {
    println!("{}", x);
}
\end{lstlisting}
\end{samepage}

The method \texttt{last} returns the last element
in the vector, and
the method \texttt{pop} removes the last element:

\begin{lstlisting}
let mut v = Vec::new();
v.push(5);
v.push(2);
println!("{}", v.last().unwrap()); // 2
v.pop();
println!("{}", v.last().unwrap()); // 5
\end{lstlisting}

The following code creates a vector with five elements:

\begin{lstlisting}
let v = vec![2,4,2,5,1];
\end{lstlisting}

Another way to create a vector is to give the number
of elements and the initial value for each element:

\begin{lstlisting}
// size 10, initial value 0
let v = vec![0; 10];
\end{lstlisting}
\begin{lstlisting}
// size 10, initial value 5
let v = vec![5; 10];
\end{lstlisting}

The internal implementation of a vector
uses an ordinary array.
If the size of the vector increases and
the array becomes too small,
a new array is allocated and all the
elements are moved to the new array.
However, this does not happen often and the
average time complexity of
\texttt{push} is $O(1)$.

\index{string}

The \texttt{String} type
is also a dynamic array that can be used almost like a vector.
In addition, there is special syntax for strings
that is not available in other data structures.
Strings can be combined using the \texttt{+} symbol.
The method \texttt{substr} is not available,
but string slices can be used to extract substrings,
and the method \texttt{find} finds the position
of the first occurrence of a substring.

The following code presents some string operations:

\begin{lstlisting}
let mut a = String::from("hatti");
let mut b = format!("{}{}", a, a);
println!("{}", b); // hattihatti
let mut chars: Vec<char> = b.chars().collect();
chars[5] = 'v';
b = chars.into_iter().collect();
println!("{}", b); // hattivatti
let c = &b[3..7];
println!("{}", c); // tiva
\end{lstlisting}

\section{Set structures}

\index{set}

A \key{set} is a data structure that
maintains a collection of elements.
The basic operations of sets are element
insertion, search and removal.

The Rust standard library contains two set
implementations:
The structure \texttt{BTreeSet} is based on a balanced
search tree (a B-tree) and its operations work in $O(\log n)$ time.
The structure \texttt{HashSet} uses hashing,
and its operations work in $O(1)$ time on average.

The choice of which set implementation to use
is often a matter of taste.
The benefit of the \texttt{BTreeSet} structure
is that it maintains the order of the elements
and provides functions that are not available
in \texttt{HashSet}.
On the other hand, \texttt{HashSet}
can be more efficient.

The following code creates a set
that contains integers,
and shows some of the operations.
The method \texttt{insert} adds an element to the set,
the method \texttt{contains} checks whether
an element is in the set,
and the method \texttt{remove} removes an element from the set.

\begin{lstlisting}
use std::collections::BTreeSet;

let mut s = BTreeSet::new();
s.insert(3);
s.insert(2);
s.insert(5);
println!("{}", s.contains(&3) as i32); // 1
println!("{}", s.contains(&4) as i32); // 0
s.remove(&3);
s.insert(4);
println!("{}", s.contains(&3) as i32); // 0
println!("{}", s.contains(&4) as i32); // 1
\end{lstlisting}

A set can be used mostly like a vector,
but it is not possible to access
the elements using the \texttt{[]} notation.
The following code creates a set,
prints the number of elements in it, and then
iterates through all the elements:
\begin{lstlisting}
use std::collections::BTreeSet;

let s = BTreeSet::from([2,5,6,8]);
println!("{}", s.len()); // 4
for x in &s {
    println!("{}", x);
}
\end{lstlisting}

An important property of sets is
that all their elements are \emph{distinct}.
Thus, the method \texttt{contains} always returns
either \texttt{false} (the element is not in the set)
or \texttt{true} (the element is in the set),
and the method \texttt{insert} never adds
an element to the set if it is
already there.
The following code illustrates this:

\begin{lstlisting}
use std::collections::BTreeSet;

let mut s = BTreeSet::new();
s.insert(5);
s.insert(5);
s.insert(5);
println!("{}", s.contains(&5)); // true
\end{lstlisting}

Rust does not have a multiset in the standard library,
but a multiset can be implemented using
\texttt{BTreeMap} where each value stores
the number of instances of the key.
The following code adds three instances
of the number 5:

\begin{lstlisting}
use std::collections::BTreeMap;

let mut s = BTreeMap::new();
*s.entry(5).or_insert(0) += 1;
*s.entry(5).or_insert(0) += 1;
*s.entry(5).or_insert(0) += 1;
println!("{}", s.get(&5).copied().unwrap_or(0)); // 3
\end{lstlisting}
The method \texttt{remove} removes
all instances of an element
from the multiset:
\begin{lstlisting}
use std::collections::BTreeMap;

let mut s = BTreeMap::new();
*s.entry(5).or_insert(0) += 3;
s.remove(&5);
println!("{}", s.get(&5).copied().unwrap_or(0)); // 0
\end{lstlisting}
Often, only one instance should be removed,
which can be done as follows:
\begin{lstlisting}
use std::collections::BTreeMap;

let mut s = BTreeMap::new();
*s.entry(5).or_insert(0) += 3;
if let Some(c) = s.get_mut(&5) {
    *c -= 1;
    if *c == 0 { s.remove(&5); }
}
println!("{}", s.get(&5).copied().unwrap_or(0)); // 2
\end{lstlisting}

\section{Map structures}

\index{map}

A \key{map} is a generalized array
that consists of key-value-pairs.
While the keys in an ordinary array are always
the consecutive integers $0,1,\ldots,n-1$,
where $n$ is the size of the array,
the keys in a map can be of any data type and
they do not have to be consecutive values.

The Rust standard library contains two map
implementations that correspond to the set
implementations: the structure
\texttt{BTreeMap} is based on a balanced
search tree (a B-tree) and accessing elements
takes $O(\log n)$ time,
while the structure
\texttt{HashMap} uses hashing
and accessing elements takes $O(1)$ time on average.

The following code creates a map
where the keys are strings and the values are integers:

\begin{lstlisting}
use std::collections::BTreeMap;

let mut m = BTreeMap::new();
m.insert("monkey".to_string(), 4);
m.insert("banana".to_string(), 3);
m.insert("harpsichord".to_string(), 9);
println!("{}", m["banana"]); // 3
\end{lstlisting}

Note that \texttt{m[key]} causes a panic
if the map does not contain the key.
If the value of a key is requested
but the map does not contain it,
the key can be added to the map with
a default value using \texttt{or\_insert}.
For example, in the following code,
the key ''aybabtu'' with value 0
is added to the map.

\begin{lstlisting}
use std::collections::BTreeMap;

let mut m = BTreeMap::new();
println!("{}", m.entry("aybabtu").or_insert(0)); // 0
\end{lstlisting}
The method \texttt{contains\_key} checks
if a key exists in a map:
\begin{lstlisting}
use std::collections::BTreeMap;

let mut m = BTreeMap::new();
m.insert("aybabtu".to_string(), 0);
if m.contains_key("aybabtu") {
    // key exists
}
\end{lstlisting}
The following code prints all the keys and values
in a map:
\begin{lstlisting}
use std::collections::BTreeMap;

let mut m = BTreeMap::new();
m.insert("monkey".to_string(), 4);
m.insert("banana".to_string(), 3);
for (k, v) in &m {
    println!("{} {}", k, v);
}
\end{lstlisting}

\section{Iterators and ranges}

\index{iterator}

Many functions in the Rust standard library
operate with slices and iterators.
An \key{iterator} is a value that allows
accessing elements in a data structure one by one.

The methods \texttt{iter} and \texttt{iter\_mut}
define a range that contains
all elements in a data structure.
The elements can be processed using a \texttt{for} loop.

The often used methods \texttt{sort}
and \texttt{reverse} process all elements
in a vector:

\begin{lstlisting}
let mut v = vec![3,1,4,1,5];
v.sort();
v.reverse();
\end{lstlisting}

The same methods can be used with an ordinary array
or a slice:

\begin{lstlisting}
let mut a = [3,1,4,1,5];
a.sort();
a.reverse();
\end{lstlisting}

Shuffling the order of elements
requires the \texttt{rand} crate:

\begin{lstlisting}
use rand::seq::SliceRandom;
use rand::thread_rng;

let mut v = vec![3,1,4,1,5];
let mut rng = thread_rng();
v.shuffle(&mut rng);
\end{lstlisting}

\subsubsection{Working with sets}

Elements of a set can be accessed
using an iterator.
The following code creates an iterator
that points to the smallest element in a set:

\begin{lstlisting}
use std::collections::BTreeSet;

let s = BTreeSet::from([3,4,6,8,12,13,14,17]);
let first = s.iter().next();
\end{lstlisting}

The element to which an iterator points
can be accessed directly.
For example, the following code prints
the first element in the set:

\begin{lstlisting}
use std::collections::BTreeSet;

let s = BTreeSet::from([3,4,6,8,12,13,14,17]);
if let Some(x) = s.iter().next() {
    println!("{}", x);
}
\end{lstlisting}

The elements of a set can be traversed
in increasing order:

\begin{lstlisting}
use std::collections::BTreeSet;

let s = BTreeSet::from([3,4,6,8,12,13,14,17]);
for x in &s {
    println!("{}", x);
}
\end{lstlisting}
The following code prints the largest element in the set:
\begin{lstlisting}
use std::collections::BTreeSet;

let s = BTreeSet::from([3,4,6,8,12,13,14,17]);
if let Some(x) = s.iter().next_back() {
    println!("{}", x);
}
\end{lstlisting}

The method \texttt{get} checks
whether the set contains an element:

\begin{lstlisting}
use std::collections::BTreeSet;

let s = BTreeSet::from([3,4,6,8,12,13,14,17]);
if s.get(&8).is_some() {
    // 8 is found
}
\end{lstlisting}

Ordered sets also support range queries:
the method \texttt{range} iterates over
elements within a given range of values.
These methods are not supported by the
\texttt{HashSet} structure which
does not maintain the order of the elements.

For example, the following code finds the element
nearest to $x$:

\begin{lstlisting}
use std::collections::BTreeSet;

let s = BTreeSet::from([3,4,6,8,12,13,14,17]);
let x = 10;
let next = s.range(x..).next();       // smallest element >= x
let prev = s.range(..x).next_back();  // largest element < x
let nearest = match (prev, next) {
    (None, Some(&a)) => a,
    (Some(&b), None) => b,
    (Some(&b), Some(&a)) => if x - b < a - x { b } else { a },
    (None, None) => unreachable!(), // the set is not empty
};
println!("{}", nearest);
\end{lstlisting}

The code assumes that the set is not empty,
and goes through all possible cases.
First, the code looks for the smallest
element whose value is at least $x$
and the largest element whose value is smaller than $x$.
If only one of them exists,
it is nearest to $x$.
Otherwise, the nearest element is the one
that is closer to $x$
(the larger one if the distances are equal).

\section{Other structures}

\subsubsection{Bitset}

\index{bitset}

A \key{bitset} is an array
whose each value is either 0 or 1.
Rust does not have a bitset in the standard library,
but for a small fixed number of bits an integer
can be used as a bitmask.
For example, the following code creates
a bitset that contains 10 elements:

\begin{lstlisting}
let mut s: u16 = 0;
s |= 1 << 1;
s |= 1 << 3;
s |= 1 << 4;
s |= 1 << 7;
println!("{}", (s >> 4) & 1); // 1
println!("{}", (s >> 5) & 1); // 0
\end{lstlisting}

The benefit of using bitsets is that
they require less memory than ordinary arrays,
because each element in a bitset only
uses one bit of memory.
For example,
if $n$ bits are stored in an \texttt{i32} array,
$32n$ bits of memory will be used,
but a corresponding bitset only requires $n$ bits of memory.
In addition, the values of a bitset
can be efficiently manipulated using
bit operators, which makes it possible to
optimize algorithms using bit sets.

The following code shows another way to create the above bitset:
\begin{lstlisting}
let mut s: u16 = 0b0010011010; // from right to left
println!("{}", (s >> 4) & 1); // 1
println!("{}", (s >> 5) & 1); // 0
\end{lstlisting}

The method \texttt{count\_ones} returns the number
of ones in the bitset:

\begin{lstlisting}
let s: u16 = 0b0010011010;
println!("{}", s.count_ones()); // 4
\end{lstlisting}

The following code shows examples of using bit operations:
\begin{lstlisting}
let a: u16 = 0b0010110110;
let b: u16 = 0b1011011000;
println!("{:010b}", a & b); // 0010010000
println!("{:010b}", a | b); // 1011111110
println!("{:010b}", a ^ b); // 1001101110
\end{lstlisting}

\subsubsection{Deque}

\index{deque}

A \key{deque} is a dynamic array
whose size can be efficiently
changed at both ends of the array.
Like a vector, a deque provides the methods
\texttt{push\_back} and \texttt{pop\_back}, but
it also includes the methods
\texttt{push\_front} and \texttt{pop\_front}
which are not available in a vector.

A deque can be used as follows:
\begin{lstlisting}
use std::collections::VecDeque;

let mut d = VecDeque::new();
d.push_back(5); // [5]
d.push_back(2); // [5,2]
d.push_front(3); // [3,5,2]
d.pop_back(); // [3,5]
d.pop_front(); // [5]
\end{lstlisting}

The internal implementation of a deque
is more complex than that of a vector,
and for this reason, a deque is slower than a vector.
Still, both adding and removing
elements take $O(1)$ time on average at both ends.

\subsubsection{Stack}

\index{stack}

A \key{stack}
is a data structure that provides two
$O(1)$ time operations:
adding an element to the top,
and removing an element from the top.
It is only possible to access the top
element of a stack.

The following code shows how a stack can be used:
\begin{lstlisting}
let mut s = Vec::new();
s.push(3);
s.push(2);
s.push(5);
println!("{}", s.last().unwrap()); // 5
s.pop();
println!("{}", s.last().unwrap()); // 2
\end{lstlisting}

\subsubsection{Queue}

\index{queue}

A \key{queue} also
provides two $O(1)$ time operations:
adding an element to the end of the queue,
and removing the first element in the queue.
It is only possible to access the first
and last element of a queue.

The following code shows how a queue can be used:
\begin{lstlisting}
use std::collections::VecDeque;

let mut q = VecDeque::new();
q.push_back(3);
q.push_back(2);
q.push_back(5);
println!("{}", q.front().unwrap()); // 3
q.pop_front();
println!("{}", q.front().unwrap()); // 2
\end{lstlisting}

\subsubsection{Priority queue}

\index{priority queue}
\index{heap}

A \key{priority queue}
maintains a set of elements.
The supported operations are insertion and,
depending on the type of the queue,
retrieval and removal of
either the minimum or maximum element.
Insertion and removal take $O(\log n)$ time,
and retrieval takes $O(1)$ time.

While an ordered set efficiently supports
all the operations of a priority queue,
the benefit of using a priority queue is
that it has smaller constant factors.
A priority queue is usually implemented using
a heap structure that is much simpler than a
balanced binary tree used in an ordered set.

\begin{samepage}
By default, the elements in a Rust
priority queue (\texttt{BinaryHeap}) are sorted in decreasing order,
and it is possible to find and remove the
largest element in the queue.
The following code illustrates this:

\begin{lstlisting}
use std::collections::BinaryHeap;

let mut q = BinaryHeap::new();
q.push(3);
q.push(5);
q.push(7);
q.push(2);
println!("{}", q.peek().unwrap()); // 7
q.pop();
println!("{}", q.peek().unwrap()); // 5
q.pop();
q.push(6);
println!("{}", q.peek().unwrap()); // 6
q.pop();
\end{lstlisting}
\end{samepage}

If we want to create a priority queue
that supports finding and removing
the smallest element,
we can do it as follows:

\begin{lstlisting}
use std::collections::BinaryHeap;
use std::cmp::Reverse;

let mut q = BinaryHeap::new();
q.push(Reverse(3));
q.push(Reverse(5));
q.push(Reverse(7));
q.push(Reverse(2));
println!("{}", q.peek().unwrap().0); // 2
q.pop();
println!("{}", q.peek().unwrap().0); // 3
q.pop();
\end{lstlisting}

\subsubsection{Policy-based data structures}

Rust does not include order-statistics
data structures in the standard library.
Such structures are available in external
crates, for example the \texttt{indexset} crate
which provides a set that can be indexed like an array.
Alternatively, a Fenwick tree (covered later in the book)
can be used to support similar queries.

\section{Comparison to sorting}

It is often possible to solve a problem
using either data structures or sorting.
Sometimes there are remarkable differences
in the actual efficiency of these approaches,
which may be hidden in their time complexities.

Let us consider a problem where
we are given two lists $A$ and $B$
that both contain $n$ elements.
Our task is to calculate the number of elements
that belong to both of the lists.
For example, for the lists
\[A = [5,2,8,9] \hspace{10px} \textrm{and} \hspace{10px} B = [3,2,9,5],\]
the answer is 3 because the numbers 2, 5
and 9 belong to both of the lists.

A straightforward solution to the problem is
to go through all pairs of elements in $O(n^2)$ time,
but next we will focus on
more efficient algorithms.

\subsubsection{Algorithm 1}

We construct a set of the elements that appear in $A$,
and after this, we iterate through the elements
of $B$ and check for each element if it
also belongs to $A$.
This is efficient because the elements of $A$
are in a set.
Using the \texttt{BTreeSet} structure,
the time complexity of the algorithm is $O(n \log n)$.

\subsubsection{Algorithm 2}

It is not necessary to maintain an ordered set,
so instead of the \texttt{BTreeSet} structure
we can also use the \texttt{HashSet} structure.
This is an easy way to make the algorithm
more efficient, because we only have to change
the underlying data structure.
The time complexity of the new algorithm is $O(n)$.

\subsubsection{Algorithm 3}

Instead of data structures, we can use sorting.
First, we sort both lists $A$ and $B$.
After this, we iterate through both the lists
at the same time and find the common elements.
The time complexity of sorting is $O(n \log n)$,
and the rest of the algorithm works in $O(n)$ time,
so the total time complexity is $O(n \log n)$.

\subsubsection{Efficiency comparison}

The following table shows how efficient
the above algorithms are when $n$ varies and
the elements of the lists are random
integers between $1 \ldots 10^9$:

\begin{center}
\begin{tabular}{rrrr}
$n$ & Algorithm 1 & Algorithm 2 & Algorithm 3 \\
\hline
$10^6$ & $1.5$ s & $0.3$ s & $0.2$ s \\
$2 \cdot 10^6$ & $3.7$ s & $0.8$ s & $0.3$ s \\
$3 \cdot 10^6$ & $5.7$ s & $1.3$ s & $0.5$ s \\
$4 \cdot 10^6$ & $7.7$ s & $1.7$ s & $0.7$ s \\
$5 \cdot 10^6$ & $10.0$ s & $2.3$ s & $0.9$ s \\
\end{tabular}
\end{center}

Algorithms 1 and 2 are equal except that
they use different set structures.
In this problem, this choice has an important effect on
the running time, because Algorithm 2
is 4–5 times faster than Algorithm 1.

However, the most efficient algorithm is Algorithm 3
which uses sorting.
It only uses half the time compared to Algorithm 2.
Interestingly, the time complexity of both
Algorithm 1 and Algorithm 3 is $O(n \log n)$,
but despite this, Algorithm 3 is ten times faster.
This can be explained by the fact that
sorting is a simple procedure and it is done
only once at the beginning of Algorithm 3,
and the rest of the algorithm works in linear time.
On the other hand,
Algorithm 1 maintains a complex balanced binary tree
during the whole algorithm.
